\label{related-work}

This section situates \sysname{} with respect to the most relevant algorithmic and system approaches for distributed training limited by the WAN, focusing on precise differences in the mechanism and objectives.

\subsection{Local Updates / Local‑SGD}
Local‑SGD reduces global communication by executing multiple local SGD steps between periodic aggregations, which lowers aggregate bandwidth but induces ``client drift''~\cite{stich2019local, stich2018sparsified, compensatefedavg}. The sysname{} function does not primarily sparsify synchronization points; rather, it uses micro-batch granular continuous gradient transmission to avoid model drift.

\subsection{Gradient Compression and Error‑Feedback}
Quantization and sparsification combined with error feedback reduce transmitted bytes while recovering convergence guarantees under certain assumptions~\cite{DBLP:conf/dac/WangWLK25, zhang2023errorfeedback, zhou2023integrating, DBLP:conf/icml/XuH22}. Those techniques principally address bias from lossy encoding; \sysname{} addresses a distinct class of temporal misalignment errors arising from early transmission of gradients and is compatible with compression + error‑feedback applied on its continuous gradient stream.

\subsection{Overlap and Pipelining}
Overlap techniques and model‑parallel pipelining hide communication by interleaving compute and communication or by partitioning layers across devices~\cite{gpipe, pipedream, zerobubble}. While such methods can reduce on‑step communication tails, cross‑DC model pipelining typically requires activation transfers that are prohibitive in WAN environments; unlike opportunistic intra‑step overlap, \sysname{} proactively moves micro‑batch gradients out‑of‑band into a sustained background flow, avoiding cross‑DC activation transfer and converting bursty synchronization into steady, low‑jitter traffic.

\subsection{Summary and Positioning}
\sysname{} draws on three related lines of work: micro-batch/overlap techniques, asynchronous and delay-compensation methods, and error-feedback approaches. It departs from these prior lines in three principal ways: (i) it moves synchronization to micro-batch granularity; (ii) it employs predictive correction targeted at the temporal mismatch between previewed gradients and subsequent local updates. The empirical section evaluates \sysname{} against representative instances from these families in terms of throughput, convergence speed, and final accuracy under diverse WAN conditions and data heterogeneity.

% \textbf{Cross-datacenter / WAN training systems.}
% Training across geographically separated datacenters is fundamentally constrained by wide-area network (WAN) properties: lower bandwidth, higher latency, and higher variability than intra-datacenter fabrics. These characteristics make collective communication a first-order bottleneck, and can turn otherwise negligible synchronization into a dominant component of iteration time. As a result, Cross-DC training systems often adopt WAN-aware orchestration (e.g., topology-aware collectives and placement) and favor parallelism layouts that minimize cross-site traffic on the critical path. In particular, using data parallelism across sites avoids shipping activations across the WAN (as would happen if pipeline boundaries span datacenters), but it still requires frequent gradient aggregation across sites, leaving a recurring synchronization barrier to optimize.

% \textbf{Communication overlap, asynchronous optimization, and Local SGD.}
% A broad line of work improves throughput by overlapping gradient communication with backpropagation, for example by launching reductions for gradient buckets as soon as they are ready. While such overlap can hide a substantial portion of communication, the final ``tail'' of an iteration often remains on the critical path because some reductions must complete before the optimizer step. Relaxing synchronization further---via asynchronous or semi-asynchronous SGD---can remove global barriers but typically introduces stale gradients and may destabilize convergence. Local-update methods (often referred to as Local SGD / federated averaging) reduce synchronization frequency by performing multiple local steps between global aggregations; however, the remaining synchronization can still be expensive under WAN conditions, and increasing the local update window can lead to model drift and degraded quality. \sysname is complementary: rather than only reducing synchronization frequency, it restructures \emph{when} gradients are synchronized so that communication can proceed continuously in the background.

% \textbf{Error feedback, delayed-gradient correction, and pipelining.}
% There is extensive literature on mitigating optimization error introduced by delayed or approximate communication, including error-feedback mechanisms used with gradient compression, correction terms for delayed gradients. These methods share the principle of tracking the discrepancy between an ``ideal'' update and the update actually applied under system constraints, then feeding that discrepancy back into future updates. \sysname follows this lineage with a predictive error-correction mechanism tailored to micro-batch-level pipelining: gradients are communicated early to enable overlap, while a lightweight correction term compensates for the mismatch induced by applying updates out of phase with subsequent gradient computation. This positions \sysname as combining fine-grained pipelined synchronization with principled correction for WAN-constrained Cross-DC training.